% Generated draft: requires mathematical and source review.
\documentclass[11pt]{article}
\usepackage[margin=25mm]{geometry}
\usepackage{fontspec}
\setmainfont{Latin Modern Roman}
\usepackage{amsmath,amssymb,mathtools,unicode-math}
\setmathfont{Latin Modern Math}
\usepackage{enumitem,needspace,xurl,hyperref}
\setlength{\parindent}{0pt}
\setlength{\parskip}{5pt}
\setlength{\emergencystretch}{3em}
\newcommand{\N}{\mathbb{N}}
\newcommand{\Z}{\mathbb{Z}}
\newcommand{\Q}{\mathbb{Q}}
\newcommand{\R}{\mathbb{R}}
\newcommand{\C}{\mathbb{C}}
\newcommand{\F}{\mathbb{F}}
\newcommand{\A}{\mathbb{A}}
\newcommand{\PP}{\mathbb{P}}
\newcommand{\E}{\mathbb{E}}
\newcommand{\eps}{\varepsilon}
\newcommand{\dd}{\,\mathrm{d}}
\begin{document}
\section*{UnsolvedMath ID: 801127}
\textbf{Source problem code:} catalogue entry
\subsection{Definitions and mathematical statement}
A type $\mathrm{II}_1$ factor is a von Neumann algebra $N$ with center $Z(N)=\C 1$, a faithful normal tracial state $\tau$ satisfying $\tau(1)=1$, and no nonzero minimal projections. Under the standard tensor-product convention, such a factor is \emph{non-prime} if there exist type $\mathrm{II}_1$ factors $P$ and $Q$ such that
\[
N\cong P\,\overline{\otimes}\,Q,
\]
where $\overline{\otimes}$ denotes the spatial von Neumann algebra tensor product.

The supplied question is:
\begin{quote}
Does every non-prime type $\mathrm{II}_1$ factor $N$ have a unique sequential commutation orbit?
\end{quote}
No separability hypothesis or restriction to ultraproduct factors is stated in the supplied problem, despite the title of its source list.

\textbf{Unresolved definition.} The supplied record does not define a sequential commutation orbit. In particular, it does not specify the objects forming the orbits or the relation that determines them. Thus the question above preserves the source wording, but a precise mathematical formulation of its uniqueness assertion requires the missing definition; no replacement definition is imposed here.

The source adds: ``See [HoudayerIoana2024] where this is proven for commutants.'' This is a supplied attribution, not an independently verified theorem statement. The record does not identify the relevant commutants, their ambient algebras, or the hypotheses of that result.
\subsection{Short English statement}
The source asks whether every type $\mathrm{II}_1$ factor that decomposes as a tensor product of two type $\mathrm{II}_1$ factors has exactly one sequential commutation orbit. The meaning of that orbit is absent from the supplied record, so the intended uniqueness assertion remains incompletely specified. The record labels the problem open and attributes a result for commutants to [HoudayerIoana2024]; neither assessment has been independently verified.
\subsection{Sources}
\begin{itemize}
\item Supplied local UnsolvedMath registry record, permanent ID $801127$, exact source problem code \texttt{catalogue entry}, titled ``Ultraproduct $\mathrm{II}_1$ Factors --- Problem 28.'' The record attributes the problem to Elayavalli, with proposed year 2025, and supplies the status \texttt{open}.
\item Supplied source-list attribution: Elayavalli, \emph{50 Open Problems: Ultraproduct $\mathrm{II}_1$ Factors} (2025), Problem 28. Supplied URL: \url{https://arxiv.org/abs/2511.20377}. The record reports access on 2026-10-16 and extraction by \texttt{arxiv\_latex}; these are supplied provenance details, not a claim of independent access or verification.
\item The statement cites [HoudayerIoana2024] and says that the assertion is proven for commutants. No full bibliographic entry, URL, theorem number, or hypotheses for this reference are supplied, so the cited result cannot be identified precisely from this record alone.
\end{itemize}
\end{document}
